% Option-A abstract for the current pre-official-campaign manuscript.
% Numeric values below are legacy reconstruction results only and must be replaced or supplemented after A3.

\begin{abstract}
Cross-chain bridge failures can arise from interactions among source-chain contracts, relay or validator logic, and destination-chain contracts, making security testing difficult to reduce to a single contract or execution domain. Existing approaches provide complementary capabilities, including cross-chain static analysis, transaction monitoring, graph-based attack detection, and dynamic fuzzing. This paper presents BridgeSentry, a semantic-guided framework for reconstructing and validating exploit traces in EVM-compatible cross-chain bridges. BridgeSentry combines (1) hybrid semantic extraction using structured program analysis with optional LLM enrichment to construct Atomic Transfer Graphs (ATGs) and candidate invariants, (2) provenance-aware retrieval of historical bridge incidents to generate threat-model-constrained adversarial action sequences, and (3) a dual-EVM execution harness with an explicit relay and semantic guidance.

\textbf{Current evaluation scope:} the retained legacy evaluation contains 12 reconstructed historical incidents: 10 EVM reconstructions, one EVM surrogate of a Solana vulnerability, and one scenario-layer-only case. Under a full-knowledge retrieval setting, all 12 target predicates are triggered, 11 cases execute target bytecode, and the mean fuzzing-only trigger time on successful legacy runs is 1.66 seconds. These quantities characterize target-aware reconstruction and do not constitute a valid-exploit rate, precision estimate, or zero-day discovery result. The current artifact does not yet contain a completed incident-disjoint, paired-patch, benign, or hard-negative campaign. We therefore treat invariant triggers as intermediate outcomes and reserve effectiveness claims for a fail-closed evaluation in which exploit evidence is derived from runtime execution, satisfies the stated adversary model, demonstrates material impact, and is rejected by a paired patched control.
\end{abstract}
